\documentclass[../../main.tex]{subfiles}

\begin{document}

\chapter{Python Basics for Machine Learning}

To develop a model represented in abstract math, we need a
programming language. If you ask any ML engineer on the
language they use the most, I can guarantee you that it's
Python (well most likely). The reason is simple: syntax and
libraries. Although all heavy computations in ML is mostly done
by C++, Python is way easier to write than C++. You don't have to
deal with segmentation faults and spend more time adding the
missing semicolon than actually debugging. Moreover, Python has
popular libraries like NumPy, Matplotlib, and PyTorch that will
make our lives easier. Therefore, we will discuss about basic
syntax, NumPy, and Matplotlib for ML in this note. Don't worry,
we will discuss about PyTorch in latter notes as it requires
ML concepts.

\section{Basic Syntax}

As always, let's get started by being more familiar with the
nature of Python. Unlike most other languages, indentation is a
vital factor in Python. Two spaces, four spaces, doesn't really
matter, let's use four spaces as it is more common and readable.
Here is the Hello World example.

\begin{jupyter}
\begin{lstlisting}[style=jin]
print("Hello World")
\end{lstlisting}

\begin{lstlisting}[style=jout]
Hello World
\end{lstlisting}
\end{jupyter}

Please note that this Jupyter Notebook can be found in
\verb|jupyter/| directory in GitHub under the root directory for
these notes. Please feel free to download and follow along!
The next example is a list of simple operations that we can
do with Python.

\begin{jupyter}
\begin{lstlisting}[style=jin]
import math
import numpy as np
import matplotlib.pyplot as plt
%matplotlib inline

print(1 + 2)       # Addition
print(3 - 9)       # Subtraction
print(4 * 3)       # Multiplication
print(5 ** 2)      # Power
print(math.log(5)) # log(a, base), default base is e
print(math.exp(2)) # e^2
print(math.e)      # Accessing constants
\end{lstlisting}

\begin{lstlisting}[style=jout]
3
-6
12
25
1.6094379124341003
7.38905609893065
2.718281828459045
\end{lstlisting}
\end{jupyter}

As you can see in Cell 2, we need module ``math'' to access
constants and functions like $e$, $\pi$, and $\log_a b$.

\begin{definition}{}{}
We can achieve the same effect with \textbf{variables}, or
a place with name that stores a value. A \textbf{function}
is a reusable block of code dedicated for specific tasks.
\end{definition}

Let's take a look at an example of how function and variable
in addition to the operations above can be used to find
the approximation of the derivative of $f(x) = x^2 + 5$
at $x = 3$.

\begin{jupyter}
\begin{lstlisting}[style=jin]
def f(x):
    return x**2 + 5

a = 3
h = 0.0001
der = (f(a + h) - f(a))/h
print(der)
\end{lstlisting}

\begin{lstlisting}[style=jout]
6.000100000012054
\end{lstlisting}
\end{jupyter}

Let's check if our answer is reasonable. Using the Power Rule,
we know that $f'(x) = 2x$. Because $f(3) = 2 \cdot 3 = 6$
and $6.000100000012054$ and $6$ are somewhat close, our result
is reasonable.

\begin{definition}{}{}
\textbf{Lists} are dynamic and ordered collection of elements
of potentially different data types. An \textbf{index} is the
position of an element in the list. The index of the first item
is $0$.
\end{definition}

Below are examples of what we can do with lists.

\begin{jupyter}
\begin{lstlisting}[style=jin]
# A list can be empty
a = []
print(a)

# We can append a value at the end
a = [1, 'hi', 3]
a.append('bye')
print(a)

# Remove an element
a.remove('hi')
print(a)

# Access an element with index
print(a[2])

# Find the number of elements
print(len(a))
\end{lstlisting}

\begin{lstlisting}[style=jout]
[]
[1, 'hi', 3, 'bye']
[1, 3, 'bye']
bye
3
\end{lstlisting}
\end{jupyter}

What do we do if we need to access multiple elements? This is
where \textbf{slicing} shines. Consider the following example.

\begin{jupyter}
\begin{lstlisting}[style=jin]
# a[i:j] to access from index 'i' to index 'j-1'
a = [1, 2, 3, 4, 5, 6, 7, 8, 9]
print(a[1:5])

# Omit 'i' or 'j' to define only the start or end values
print(a[:6])
print(a[3:])
\end{lstlisting}

\begin{lstlisting}[style=jout]
[2, 3, 4, 5]
[1, 2, 3, 4, 5, 6]
[4, 5, 6, 7, 8, 9]
\end{lstlisting}
\end{jupyter}

Now that we know basic operations, lists, and how we can use
functions, let's discuss about conditions and loops.

\begin{definition}{}{}
\textbf{Conditional statements} are statements that performs
certain tasks under specific conditions. \textbf{Loops} are
structure that repeats a block of code until they are explicitly
told to end. One more definition that will make our lives easier is
\textbf{increment}, which is an assignment operator that augments
certain values.
\end{definition}

Here is an example of conditional statements.

\begin{jupyter}
\begin{lstlisting}[style=jin]
# If, Else, Elif (else if)
banana_num = 10
if banana_num < 10:
    print("There are less than ten bananas")
elif banana_num == 10:
    print("There are ten bananas")
else:
    print("There are more than ten bananas")
\end{lstlisting}

\begin{lstlisting}[style=jout]
There are ten bananas
\end{lstlisting}
\end{jupyter}

Here, because ``$=$'' is reserved, we use ``$==$'' for ``equal to.''
We could also do less than or equal to, greater than or equal to,
and not equal to with $<=$, $>=$, and $!=$ respectively.

Regarding loops, ``for loop'' and ``while loop'' are two primary
methods. Below are examples of how both loops can be used.

\begin{jupyter}
\begin{lstlisting}[style=jin]
# While loop
love_banana = True
banana_num = 0

while love_banana and banana_num < 5:
    banana_num += 1 # Increment
    print(banana_num)

print()

# For loop
my_banana = ['yellow', 'green', 'mushy', 'sweet']

for i in range(len(my_banana)):
    print(my_banana[i])
\end{lstlisting}

\begin{lstlisting}[style=jout]
1
2
3
4
5

yellow
green
mushy
sweet
\end{lstlisting}
\end{jupyter}

First, \verb|love_banana = True| is a Boolean value.
In the while loop, what it does is while \verb|love_banana| is True,
we will increase the value \verb|banana_num| by 1 and print them.
Notice that the increment with \verb|banana_num += 1| is a shorter
form of writing \verb|banana_num =| \\ \verb|banana_num + 1|.
Moreover, we use \verb|and| to make the loop continue if
\verb|love_banana| is true and \verb|banana_num| is less than five.

In the second loop, we listed all elements in the list
\verb|my_banana|. The loop \verb|for i in range()| tells Python to
go through the code \\ \verb|print(my_banana[i])| from $i = 0$ to
$\text{range(len(my\_banana))} - 1$.

We discussed very basic Python syntax that we must know to
understand the following notes. Let's discuss core libraries
NumPy and Matplotlib now, and save PyTorch for later.

\end{document}

